\section{L1缓存命中率为零的原因分析}

\subsection*{摘要}
本章继续探讨矩阵加法应用程序的性能分析，重点聚焦于 L1 和 L2 缓存命中率。通过回顾 NVIDIA GPU（以 Ampere 架构为例）的内存层次结构，深入解析为何在矩阵加法中 L1 缓存命中率为零，并揭示 NVIDIA Nsight Compute 工具在统计 L2 缓存命中率时将所有写入操作视为"命中"的特殊机制。最后通过向量加法案例验证了这一结论，强调了在性能评估中区分读取命中率与总命中率的重要性。

\subsection{引言：矩阵加法回顾}
本章将继续讨论之前的应用程序案例，即两个矩阵相加并将结果存储在第三个矩阵中。具体而言，将矩阵 A 的所有元素与矩阵 B 的对应元素相加，结果存入矩阵 C 的对应位置（例如 $4+7=11$）。在上一章中，已经讨论了矩阵相加的原理和概念。

本章的主要目的是进行性能分析，重点关注 L1 和 L2 缓存的命中率。

\subsection{NVIDIA GPU 内存层次结构回顾}
为了理解缓存行为，需要回顾 NVIDIA GPU 的硬件实现。以下讲解基于 NVIDIA 白皮书中的内存层次结构图（以 Volta 或 Ampere 架构为例）。


\subsection{全局内存与 HBM}
从最底层的全局内存（设备内存）开始，现代高端 GPU 通常采用 HBM（High Bandwidth Memory，高带宽内存）
。相较于 DDR3/4/5 等传统内存，HBM 以其极高的带宽著称。以 A100 GPU 为例，其显存总线宽度超过 5000 位（远超 CPU 端常见的 64 位），带宽可达每秒 1.5 至 1.9 TB。这是一个极其庞大的数据传输能力。


\subsection{L2 缓存}
L2 缓存被划分为多个分区（Slice），在所有 SM 之间共享。HBM 通过内存控制器连接到 L2 缓存。无论请求
来自哪个 SM，数据都可以通过 L2 缓存进行发送或接收。对于 HBM2e 等规格，L2 缓存同样包含多个物理分区，数据可根据 SM 的
访问请求在分区间移动。

\subsection{L1 数据缓存与共享内存}
深入到单个 SM 内部，L1 数据缓存与共享内存（Shared Memory）共用同一个物理存储单元。在 Ampere 架构中，该单元的最大容量为 192 KB。开发者可以配置应用程序，将此容量按需分配给共享内存和 L1 数据缓存（例如：100 KB 用于共享内存，92 KB 用于 L1 缓存）。

需要注意的是：
\begin{itemize}
    \item L1 缓存仅在单个 SM 内部共享，不在不同 SM 之间共享。
    \item 共享内存同理，仅局限于 SM 内部。
\end{itemize}

\subsection{数据访问流程与 L1 命中率分析}

\subsection{数据加载路径}
当应用程序运行时，数据通常初始存储于 HBM 中。当 SM 内的计算核心（ALU、Tensor Core 等）需要全局内存数据时，流程如下：
\begin{enumerate}
    \item 核心向加载/存储单元（Load/Store Unit）发送请求。
    \item 加载/存储单元首先查询 L1 缓存。若命中，数据直接返回核心；若未命中，请求转发至 L2 缓存。
    \item 若 L2 缓存命中，数据返回 SM；若仍未命中，则向 HBM 发起请求。
\end{enumerate}

\subsection{为何矩阵加法的 L1 命中率为零？}
在矩阵加法 $C = A + B$ 中，每个元素仅被访问一次：读取 A 的元素，读取 B 的元素，计算后写入 C。由于没有数据复用（Reuse），同一个元素不会被再次读取。因此，理论上该应用的 L1 缓存命中率应为 0。这与在 Nsight Compute 中观察到的结果一致。

相比之下，在矩阵乘法中，矩阵 A 的一行会与矩阵 B 的多列反复相乘，存在显著的数据复用，因此会产生大于零的 L1/L2 命中率。

\textbf{注意：} 块大小（Block Size）的配置可能会影响实际的缓存行为，但基本原理不变。

\subsection{L1 缓存技术细节与合并访问}
关于 L1 缓存的关键参数：
\begin{itemize}
    \item \textbf{最大容量：} 192 KB（可配置）。
    \item \textbf{缓存行大小：} 128 Bytes。
\end{itemize}

\subsection{线程束（Warp）与内存合并}
每个线程束由 32 个线程组成。当一个线程束执行加载指令时，32 个线程在同一周期内发出请求。加载/存储单元会根据地址连续性将
这些请求合并（Coalesce）。

\textbf{示例：} 假设 32 个线程各加载一个 4 字节的浮点数，且内存地址连续。总数据量为 $32 \times 4 = 128$ 字节，恰好等于一个缓存行的大小。此时，加载/存储单元会将 32 个独立请求合并为单次 L1 缓存行请求，极大提升效率。

\subsection{L2 缓存命中率的"秘密"：写入操作被视为命中}

\subsection{观测到的异常现象}
在矩阵加法测试中，虽然 L1 命中率为 0，但 L2 总命中率却显示约为 33\%。这与理论预期不符。

\subsection{原因解析}
这是 NVIDIA Nsight Compute 工具的一个特殊统计机制：\textbf{L2 缓存中的所有写入（Store）操作均被计为"命中"，即使实际并未命中缓存。}

在矩阵加法中，每个元素对应三次操作：
\begin{itemize}
    \item 加载 A（Read）：未命中
    \item 加载 B（Read）：未命中
    \item 存储 C（Write）：\textbf{被工具计为命中}
\end{itemize}
因此，命中率计算公式为：\[\text{L2 Hit Rate} = \frac{\text{Write Ops}}
{\text{Read Ops}+ \text{Write Ops}} = \frac{100}
{100 + 100 + 100} \approx 33.3\%\]

\subsection{验证实验}
通过向量加法程序验证了这一结论：
\begin{lstlisting}[language=bash, basicstyle=\ttfamily\small]
使用Nsight Compute运行分析
ncu --metrics l1tex__hit_rate,lts__hit_rate ./vector_add
\end{lstlisting}

结果显示：
\begin{itemize}
    \item L1 命中率：0\%
    \item L2 总命中率：33.33\%
    \item L2 读取命中率（lts\_read\_hit\_rate）：$\approx 0\%$
    \item L2 写入命中率（lts\_write\_hit\_rate）：100\%
\end{itemize}

\subsection{关键结论与建议}
\begin{enumerate}
    \item 在矩阵加法等无数据复用的场景中，L1 缓存命中率确实为 0。
    \item NVIDIA Nsight Compute 将 L2 写入操作强制计为命中，导致 L2 总命中率虚高。
    \item \textbf{性能评估建议：} 在评估 L2 缓存性能时，应仅关注\textbf{读取命中率}（Read Hit Rate），而非总命中率，以获得公平准确的评估。此规则不适用于 L1 缓存。
\end{enumerate}

\subsection*{结语}
本章到此结束。矩阵加法中缓存命中率的成因及工具统计特性，如有疑问，可发送邮件或在 NVIDIA 开发者论坛中搜索相关讨论，那里包含大量宝贵信息。

